% cap01/1.1_diagnostico.tex — Diagnóstico y enunciado del problema.
\section{Diagnóstico y Enunciado del Problema}

Las organizaciones de voluntariado y del tercer sector enfrentan el reto mayúsculo de gestionar procesos organizacionales que, por su propia naturaleza descentralizada y altruista, suelen estar altamente fragmentados. Estos procesos abarcan un amplio espectro de la gobernanza: desde la incorporación y retención de voluntarios (\textit{onboarding}), la planificación operativa y logística, hasta la asignación y el seguimiento de tareas especializadas. Más crítico aún es el resguardo, el procesamiento y el archivo de información clínica, financiera y personal extremadamente sensible que fluye a través de estas entidades operativas. Según el Registro Nacional de Organizaciones No Gubernamentales de Desarrollo de la Agencia Peruana de Cooperación Internacional (APCI), existen más de un millar de organizaciones formales operando activamente en el país \citep{apci2024}. Asimismo, el Instituto Nacional de Estadística e Informática (INEI) reporta que las asociaciones civiles y organizaciones sin fines de lucro representan una porción significativa y en constante crecimiento de la demografía empresarial formal \citep{inei2024}.

A pesar de esta evidente magnitud institucional y de la creciente demanda por intervención social, estas organizaciones, en su gran mayoría, padecen un evidente rezago tecnológico. Cuando los procesos operativos de misión crítica se soportan en herramientas ofimáticas aisladas u hojas de cálculo compartidas de manera precaria, surgen tres problemas estructurales y sistémicos.

Primero, la ausencia de trazabilidad. La carencia de un sistema unificado y con capacidad de registro transaccional dificulta reconstruir de forma auditable el historial de manipulación de los datos (quién realizó qué acción, cuándo y sobre qué dato), un aspecto no negociable para las exigencias de transparencia, rendición de cuentas financieras ante cooperantes internacionales y certificaciones gubernamentales.

Segundo, el problema del \textit{data bleeding} y la dispersión de datos vulnerables. La manipulación de fichas médicas de beneficiarios, documentos de identidad de voluntarios y datos bancarios eleva drásticamente el riesgo de exposición accidental o fuga deliberada de información. Este riesgo se exacerba cuando no existe un mecanismo centralizado para hacer cumplir el \textit{Principio de Mínimo Privilegio}, lo que permite que operadores de bajo rango visualicen información transversal.

Tercero, el problema de la Escalabilidad Económica. La falta de plataformas unificadas o estándares en la industria del software diseñados exclusivamente para el modelo de negocio de una ONG obliga a estas entidades a adquirir o alquilar infraestructura de TI individual, asumiendo licencias, servidores y costos de mantenimiento que diezman su liquidez. Esta multiplicidad de esfuerzos aislados limita severamente su capacidad para escalar sus programas sociales, ya que destinan valiosos recursos de donaciones a mantener una infraestructura tecnológica fragmentada.

\textit{Democra} aborda directa y estructuralmente esta problemática mediante el diseño de una plataforma única tipo \textit{Software as a Service} (SaaS) \textit{multi-tenant}, que consolida y centraliza los múltiples dominios operativos (recursos humanos, clínica, finanzas) bajo un modelo estricto de seguridad compartimentada. A través del modelo \textit{multi-tenant} o de <<arquitectura de inquilinos>>, cada organización (\textit{tenant}) opera de manera aislada e independiente sobre una macroinfraestructura compartida, aprovechando los beneficios económicos de la nube sin sacrificar la confidencialidad gracias a una separación matemática garantizada a nivel estructural en la base de datos relacional.
